\documentclass{article}
\usepackage[margin=0.1in]{geometry}

\usepackage{tikz}
\usetikzlibrary{arrows.meta,positioning,shapes.callouts}

\begin{document}

\begin{center}
\textbf{\Large Test 1 — Database Management Systems}

\textbf{CSC 3300 \quad Spring 2026}
\end{center}

\section*{Part I — True/False}
\begin{enumerate}
  \item \textbf{(2 points)} \textbf{T/F:} In the University database, there can be a \texttt{classroom} in which no course sections were ever taught.
  \item \textbf{(2 points)} \textbf{T/F:} In the University database, each instructor can serve as an advisor to at most one student.
  \item \textbf{(2 points)} \textbf{T/F:} SQL prevents any insert to the database that violates an integrity constraint.
  \item \textbf{(2 points)} \textbf{T/F:} A student can be registered for more than one section of the same course in a given semester and year.
  \item \textbf{(2 points)} \textbf{T/F:} It can be the case that in the database there are only the following two tuples:

  \medskip
  \noindent\texttt{student}
  \begin{center}
  \begin{tabular}{|c|c|c|c|}
  \hline
  \textbf{ID} & \textbf{name} & \textbf{dept\_name} & \textbf{tot\_cred} \\\hline
  0 & Green & Comp.\ Sci. & 3 \\\hline
  \end{tabular}
  \end{center}

  \noindent\texttt{takes}
  \begin{center}
  \begin{tabular}{|c|c|c|c|c|c|}
  \hline
  \textbf{ID} & \textbf{course\_id} & \textbf{sec\_id} & \textbf{semester} & \textbf{year} & \textbf{grade} \\\hline
  0 & 3300 & 002 & Spring & 2023 & A \\\hline
  \end{tabular}
  \end{center}

  \item \textbf{(2 points)} \textbf{T/F:} In the University database, a \texttt{course\_id} can be an empty string (\texttt{''}).
  \item \textbf{(2 points)} \textbf{T/F:} In the University database, courses with the same name (i.e., the same \texttt{title}) can only exist if they are offered by distinct departments.
  \item \textbf{(2 points)} \textbf{T/F:} In the University database, a \texttt{classroom} can have 10000 seats.
  \item \textbf{(2 points)} \textbf{T/F:} In the University database, the number of tuples in the \texttt{advisor} table can exceed the number of tuples in the \texttt{student} table.
  \item \textbf{(2 points)} \textbf{T/F:} The University database maintains a record of all past advisors for each student.
  \item \textbf{(2 points)} \textbf{T/F:} The University database allows multiple \texttt{classroom}s in the same building to share the same \texttt{room\_number}.
  \item \textbf{(2 points)} \textbf{T/F:} A tuple is inserted into the \texttt{advisor} table only if the value of its \texttt{i\_ID} attribute equals the value of the \texttt{ID} attribute of some tuple in the \texttt{instructor} table.
  \item \textbf{(2 points)} \textbf{T/F:} The following two queries are equivalent, that is, they produce the same result, for any instance of the University Database:
  \begin{enumerate}
    \item select distinct course\_id from course
    \item select course\_id from course
  \end{enumerate}
\end{enumerate}

\section*{Part II — Multiple choice (choose one)}
\begin{enumerate}
  \item \textbf{(3 points)} With only 6 tuples in the \texttt{instructor} table, 2 tuples in the \texttt{section} table, and no tuples in the \texttt{teaches} table, what is the maximum number of tuples that can be added to \texttt{teaches} without violating any \emph{foreign key} constraints?
  \begin{enumerate}
    \item[(A)] 2
    \item[(B)] 4
    \item[(C)] 6
    \item[(D)] 12
  \end{enumerate}
  \item \textbf{(3 points)} Assume that there are no tuples in any of the relation instances of the University database. What is the minimal number of tuples that need to be added to the database \emph{before} one tuple can be added to the table \texttt{section}?
  \begin{enumerate}
    \item[(A)] 5
    \item[(B)] 0
    \item[(C)] 1
    \item[(D)] 2
  \end{enumerate}
\end{enumerate}

\section*{Part III — Multiple select (select all that apply)}
\begin{enumerate}

  \item \textbf{(4 points)} In which order can the following commands be invoked on an empty University database instance?

\begin{verbatim}
0  insert into instructor values ('98345', 'Kim', 'Elec. Eng.', '80000')
1  insert into classroom values ('Taylor', '3128', '70')
2  insert into department values ('Elec. Eng.', 'Taylor', '85000')
3  insert into section values ('EE-181', '1', 'Spring', '2009', 'Taylor', '3128', 'C')
4  insert into student values ('70557', 'Snow', 'Physics', '0')
5  insert into course values ('EE-181', 'Intro. to Digital Systems', 'Elec. Eng.', '3')
6  insert into department values ('Physics', 'Watson', '70000')
7  insert into time_slot values ('C', 'M', '11', '0', '11', '50')
8  insert into takes values ('70557', 'EE-181', '1', 'Spring', '2009', 'B')
9  insert into teaches values ('98345', 'EE-181', '1', 'Spring', '2009')
\end{verbatim}

  \begin{enumerate}
    \item[(A)] 6, 2, 1, 5, 0, 4, 7, 3, 8, 9
    \item[(B)] 2, 6, 5, 1, 0, 4, 7, 3, 9, 8
    \item[(C)] 2, 0, 5, 4, 3, 1, 7, 6, 8, 9
    \item[(D)] 6, 2, 0, 4, 5, 1, 7, 3, 8, 9
  \end{enumerate}




  \item \textbf{(4 points)} Which of the following SQL queries retrieve the names of departments
along with the titles of courses in the catalog of that department?
%answers A and D are correct

The result must be:
\begin{itemize}
    \item Ordered by department name (ascending)
    \item Then by course title (ascending)
\end{itemize}

\textbf{Hint:} You may find it useful to add the following course to the database
initialized with the small relations insert file provided with this exam
(\texttt{smallRelationsInsertFile\_MySQL.sql}):

\begin{verbatim}
insert into course values ('3300', 'Database Management Systems', null, '3');
\end{verbatim}

  \begin{enumerate}

    \item[(A)] 
\begin{verbatim}
select department.dept_name, title
from department, course
where department.dept_name = course.dept_name
order by dept_name, title;
\end{verbatim}

    \item[(B)] 
\begin{verbatim}
select dept_name, title
from course 
order by dept_name, title;
\end{verbatim}

    \item[(C)] 
\begin{verbatim}
select department.dept_name, title
from department, course
order by dept_name, title;
\end{verbatim}

    \item[(D)] 
\begin{verbatim}
select dept_name, title
from course
where dept_name is not null
order by dept_name, title;

\end{verbatim}

  \end{enumerate}

\end{enumerate}

\end{document}
